\documentclass[11pt,a4paper]{article}
\usepackage[T1]{fontenc}
\usepackage[utf8]{inputenc}
\usepackage{lmodern}
\usepackage{amsmath,amssymb}
\usepackage[margin=25mm]{geometry}
\usepackage[hidelinks]{hyperref}
\hypersetup{pdftitle={Compact U(1) under refinement: an explicit total-variation bound in three dimensions, and },pdfauthor={navstokgap research working notes}}
\newcommand{\tightlist}{\setlength{\itemsep}{0pt}\setlength{\parskip}{0pt}}
\setlength{\emergencystretch}{3em}
\setcounter{secnumdepth}{0}
\begin{document}
\hypertarget{compact-u1-under-refinement-an-explicit-total-variation-bound-in-three-dimensions-and-why-it-fails-in-four}{%
\section{Compact U(1) under refinement: an explicit total-variation
bound in three dimensions, and why it fails in
four}\label{compact-u1-under-refinement-an-explicit-total-variation-bound-in-three-dimensions-and-why-it-fails-in-four}}

\textbf{Result, 2026-09-27.} For the heat-kernel (Villain) \(U(1)\)
lattice gauge theory on the periodic lattice \((\mathbb Z/N)^D\) of
spacing \(a\), physical side \(L=Na\) and plaquette heat time \(t\), the
lattice measure splits exactly into a monopole-free part (the free
photon together with the physical flux sectors of the torus) and a gas
of lattice monopoles with Coulomb weights (Proposition 1; the duality of
Banks, Myerson and Kogut). A bound on the lattice Laplacian then gives:

\begin{itemize}
\tightlist
\item
  \textbf{\(D=3\) (Theorem 2).} With \(t=\lambda_3a\le1\),
  \[\bigl\|\mu_{a}-\mu_{a}^{\,0}\bigr\|_{\rm TV}\ \le\
  \frac{2(L/a)^3\,e^{-\pi^2/(6\lambda_3a)}}{1-e^{-\pi^2/(2\lambda_3a)}},\]
  where \(\mu_a^0\) is the monopole-free measure. At fixed \(\lambda_3\)
  and \(L\) the monopole free-energy density per physical volume lies
  between \(6a^{-3}e^{-2\pi^2/(3\lambda_3a)}(1+o(1))\) and
  \(a^{-3}\cdot2e^{-\pi^2/(6\lambda_3a)}/(1-e^{-\pi^2/(2\lambda_3a)})\),
  and everything compact vanishes faster than any power of \(a\).
  Göpfert and Mack's fixed-Debye-mass limit lies on a different
  trajectory, \(\lambda_3a\simeq c_0/(2\log(1/a))\) with
  \(c_0\approx4.99\) the monopole exponent, where the monopole density
  per physical volume diverges. The flux sectors carry the weight
  \(e^{-2\pi^2|m|^2/(\lambda_3L)}\), independent of \(a\).
\item
  \textbf{\(D=4\) (Theorem 3, and a recorded failure).} With \(t=g^2\),
  the same decomposition holds with flux-sector weight
  \(e^{-2\pi^2|m|^2/g^2}\), independent of both \(a\) and \(L\), and the
  monopole-loop free energy per lattice cell is at most
  \(8e^{-\pi^2/(8g^2)}/(1-e^{-3\pi^2/(8g^2)})\). The total-variation
  bound is then of order \((L/a)^4e^{-\pi^2/(8g^2)}\), useless as
  \(a\to0\) at fixed \(g\): the criterion that closes the
  three-dimensional case gives nothing here, because the heat time does
  not decrease. That the route fails, and not only the bound, needs an
  extensive lower bound on the loop free energy (a dilute-gas estimate,
  expected and not given here). The four-dimensional continuum statement
  (Driver's renormalized free field) needs observables and a coupling
  renormalization, which a partition-function bound cannot supply.
\end{itemize}

This is the first step of the \(D=3\) programme of STATE carried out at
the level of the measure, with explicit \(a\), \(L\) and coupling, for
the smallest compact group. It sharpens Theorem 5 of the
\href{series-parallel-gauge-refinement.md}{series/parallel note} (one
step, rate \(e^{-\pi^2/(8t)}\)) into a bound on the whole lattice
measure with rate \(e^{-\pi^2/(6t)}\), and it puts numbers on the
error-budget row of the \href{zero-spacing-any-action.md}{zero-spacing
note}. The duality is established
(\href{https://doi.org/10.1016/0550-3213(77)90129-8}{Banks, Myerson and
Kogut 1977}, metadata) and the continuum limit at fixed coupling is
Gross's theorem (\href{https://doi.org/10.1007/BF01210842}{Gross 1983},
abstract as indexed); the explicit total-variation bound and the
flux-sector bookkeeping are elementary consequences, with no novelty
claimed.

\hypertarget{the-exact-decomposition}{%
\subsection{1. The exact decomposition}\label{the-exact-decomposition}}

Let \(E\), \(P\) and \(C_3\) be the links, plaquettes and 3-cells of the
periodic hypercubic lattice \((\mathbb Z/N)^D\), \(d\) the coboundary on
cochains, and \(\Delta=dd^{\sf T}+d^{\sf T}d\) the Hodge Laplacian. The
heat-kernel measure on link angles \(\theta\in[0,2\pi)^E\) is

\[\mu_a(d\theta)=\frac1Z\prod_{p\in P}k_t\bigl((d\theta)_p\bigr)\,
\frac{d\theta}{(2\pi)^{|E|}},\qquad
k_t(\phi)=\sum_{n\in\mathbb Z}e^{-tn^2/2}e^{in\phi}
=\sqrt{\frac{2\pi}t}\sum_{n\in\mathbb Z}e^{-(\phi+2\pi n)^2/(2t)}.\]

Write
\(\mathbb R^P=\operatorname{im}d\oplus\operatorname{im}d^{\sf T}\oplus\mathcal H\)
(exact, coexact and harmonic 2-cochains; \(\mathcal H\) is spanned by
the cochains constant on each of the \(\binom D2\) plaquette
orientations). For \(n\in\mathbb Z^P\) let \(q=dn\in\mathbb Z^{C_3}\)
(the monopole charges) and \(n_{\mathcal H}\) the harmonic projection of
\(n\).

\textbf{Proposition 1 (monopole decomposition).} \(\mu_a\) is a mixture

\[\mu_a=\frac1Z\sum_{[n]\in\mathbb Z^P/d\mathbb Z^E}
w([n])\,\nu_{[n]},\qquad
w([n])=\exp\Bigl\{-\frac{2\pi^2}t\Bigl(q^{\sf T}(dd^{\sf T})^+q
+|n_{\mathcal H}|^2\Bigr)\Bigr\},\]

of probability measures \(\nu_{[n]}\) on link angles, where
\((dd^{\sf T})^+\) is the pseudo-inverse on \(\operatorname{im}d\) and
\(|\cdot|\) the Euclidean norm on cochains. The classes \([n]\) with
\(q=0\) make up the monopole-free measure \(\mu_a^0\).

\emph{Proof.} Insert the Poisson form of \(k_t\) and exchange sum and
integral. The substitution \(n\mapsto n+d\ell\),
\(\theta\mapsto\theta-2\pi\ell\) (\(\ell\in\mathbb Z^E\)) groups the sum
over \(n\) into classes and extends the angle integral to
\(\mathbb R^E\) modulo \(2\pi\) times the integer closed 1-cochains, a
lattice of full rank in the closed 1-cochains; the remaining integrand
depends on \(\theta\) only through \(\omega=d\theta\). So the class
\([n]\) contributes a constant times
\(\int_{\operatorname{im}d}e^{-|\omega+2\pi n|^2/(2t)}d\omega\).
Decompose \(n=n_{\rm ex}+n_{\rm co}+n_{\mathcal H}\) orthogonally. The
exact part is absorbed by translating \(\omega\), and
\(|n_{\rm co}|^2=(dn)^{\sf T}(dd^{\sf T})^+(dn)\) because
\(n_{\rm co}=d^{\sf T}(dd^{\sf T})^+dn\). The remaining Gaussian
integral is the same for every class. \(\square\)

The classes are labelled by \(q\in d\mathbb Z^P\) together with an
integer flux vector \(m\in\mathbb Z^{\binom D2}\) through the coordinate
2-tori (the integral cohomology of the torus is free). For fixed \(q\)
the harmonic part is \(n_{\mathcal H}=h(m)+y(q)\) with \(h(m)\) the
constant cochain of flux \(m\) and \(y(q)\) an offset fixed by \(q\). A
cochain of flux \(m_{\mu\nu}\) through each \((\mu\nu)\) torus has value
\(m_{\mu\nu}/N^2\) on each of the \(N^D\) plaquettes of that
orientation, so

\[|h(m)|^2=N^{D-4}\sum_{\mu<\nu}m_{\mu\nu}^2 .\]

\textbf{Two elementary facts.} (i) Each Fourier mode of \(\Delta\) on
the hypercubic lattice has eigenvalue \(\sum_\mu4\sin^2(k_\mu/2)\le4D\),
so \(\|dd^{\sf T}\|\le4D\) and \(q^{\sf T}(dd^{\sf T})^+q\ge|q|^2/(4D)\)
on \(\operatorname{im}d\). (ii) For \(y\) fixed, the theta-type sum
\(\sum_me^{-\alpha|h(m)+y|^2}\) lies between
\(e^{-\alpha|y|^2}\sum_me^{-\alpha|h(m)|^2}\) (pair \(m\) with \(-m\))
and \(\sum_me^{-\alpha|h(m)|^2}\) (its Fourier coefficients are
positive).

\textbf{The elementary monopole pair or loop.} For a single plaquette
\(p\), the charge \(q=d\,e_p\) is a nearest monopole pair (\(D=3\)) or
an elementary monopole loop (\(D=4\)), and
\(q^{\sf T}(dd^{\sf T})^+q=|P_{\rm co}e_p|^2\). By translation
invariance each plaquette has the same projections, so
\(|P_{\rm ex}e_p|^2={\rm rank}\,d_1/|P|\) and
\(|P_{\mathcal H}e_p|^2=\binom D2/|P|\). With
\({\rm rank}\,d_1=|E|-(N^D-1)-D\):

\[D=3:\ |P_{\rm co}e_p|^2=\frac13\Bigl(1-\frac1{N^3}\Bigr),\qquad
D=4:\ |P_{\rm co}e_p|^2=\frac12\Bigl(1-\frac1{N^4}\Bigr).\]

For \(D=3\) this agrees with the lattice Green's function:
\(G(0)-G(e_1)=\frac16(1-N^{-3})\), so a nearest pair has energy
\(2(G(0)-G(e_1))\).

\hypertarget{three-dimensions-the-theorem}{%
\subsection{2. Three dimensions: the
theorem}\label{three-dimensions-the-theorem}}

\textbf{Theorem 2.} Let \(D=3\), \(N=L/a\), \(t=\lambda_3a\le1\), and
\(R=Z/Z^0\) the ratio of the full partition function to its
monopole-free part. Then

\[1+6N^3\exp\Bigl\{-\frac{2\pi^2}{3\lambda_3a}
-\frac{2\pi^2a^2}{\lambda_3L^3}\Bigr\}\ \le\ R\ \le\
\exp\Bigl\{\frac{2N^3e^{-\pi^2/(6\lambda_3a)}}{1-e^{-\pi^2/(2\lambda_3a)}}\Bigr\},\]

and

\[\|\mu_a-\mu_a^0\|_{\rm TV}\le1-\frac1R\le
\frac{2(L/a)^3e^{-\pi^2/(6\lambda_3a)}}{1-e^{-\pi^2/(2\lambda_3a)}}.\]

The monopole-free part is the free photon on the torus with the flux
sectors weighted by \(\exp\{-2\pi^2|m|^2/(\lambda_3L)\}\).

\emph{Proof.} The flux weight is \(|h(m)|^2=N^{-1}|m|^2\) and
\(tN=\lambda_3L\). Upper bound: by fact (ii) the sum over \(m\) at fixed
\(q\) is at most its value at \(q=0\), so
\(R\le\sum_{q\in d\mathbb Z^P} e^{-(2\pi^2/t)q^{\sf T}(dd^{\sf T})^+q}\).
By fact (i) with \(4D=12\) each term is at most
\(e^{-(\pi^2/6t)|q|^2}\). Enlarge the sum to all of \(\mathbb Z^{C_3}\),
\(|C_3|=N^3\), and use
\(\sum_{k\in\mathbb Z}e^{-\alpha k^2}\le1+2e^{-\alpha}/(1-e^{-3\alpha})\)
(from \(k^2-1\ge3(k-1)\)) with \(\alpha=\pi^2/(6t)\), then
\(\log(1+x)\le x\). Lower bound: keep \(q=0\) and the \(6N^3\) nearest
pairs \(q=\pm d\,e_p\) (\(3N^3\) plaquettes, two signs). Each pair has
Coulomb energy at most \(\frac13\) and harmonic offset
\(|y|^2=|P_{\mathcal H}e_p|^2=N^{-3}\), and fact (ii) gives the stated
factor. Total variation: \(\mu_a\) is the mixture
\(R^{-1}\mu_a^0+(1-R^{-1})\mu_a^{\rm rest}\), so
\(\|\mu_a-\mu_a^0\|_{\rm TV}\le1-R^{-1}\le\log R\). \(\square\)

\textbf{Reading.} In a box of fixed physical side, the lattice measure
is within total variation \(O((L/a)^3e^{-\pi^2/(6\lambda_3a)})\) of the
monopole-free measure. At fixed \(\lambda_3\) this tends to zero faster
than any power of \(a\). By the summable-error criterion of the
\href{refinement-composition-and-limit.md}{refinement note}, Proposition
4, the consistency error of the compact theory is bounded by that of the
free photon with flux sectors plus a summable total-variation term; the
Gaussian comparison itself is not summable in total variation (see after
Corollary 2\('\)) and must be made through smooth observables, and
Gross's theorem is the continuum statement. At fixed \(\lambda_3\) and
\(L\) the monopole free-energy density \(L^{-3}\log R\) lies between
\(6a^{-3}e^{-2\pi^2/(3\lambda_3a)}(1+o(1))\) and the upper bound of the
summary; the lower bound is a single-defect bound, extensive only while
\(N^3e^{-2\pi^2/(3t)}\ll1\), and an extensive lower bound elsewhere
needs a dilute-gas estimate. The monopole exponent \(c_0\) lies between
the bounds' \(\pi^2/6\) and \(2\pi^2/3\) and equals
\(2\pi^2G(0)\approx4.99\) (\(G(0)\approx0.2527\)) on the infinite
lattice, single monopoles dominating. The density per physical volume
stays finite along \(\lambda_3a\simeq c_0/(3\log(1/a))\), where the
Debye mass \(m_D^2\propto\rho/\lambda_3\) tends to zero. Göpfert and
Mack's fixed-Debye-mass limit
(\href{https://doi.org/10.1007/BF01961240}{Göpfert--Mack 1982}, abstract
as indexed) is the trajectory \(\lambda_3a\simeq c_0/(2\log(1/a))\), on
which the density per physical volume diverges like \(1/(a\log(1/a))\)
and \(\lambda_3\to\infty\), consistent with their string tension over
\(m_D^2\) diverging. (Corrected 2026-09-27 after a Fable review: the
first version placed the Göpfert--Mack limit on the constant-density
trajectory.)

\textbf{Corollary 2\('\) (iteration reduces to the free field).} Let
\(a_n=2^{-n}a_0\) and let \(p_{n0}\) be the composite of the \(3n\)
directional halvings from spacing \(a_n\) to \(a_0\) (multiplying fine
links). Then

\[\bigl\|p_{n0*}\mu_{a_n}-p_{n0*}\mu_{a_n}^0\bigr\|_{\rm TV}\ \le\
\frac{2(L/a_n)^3e^{-\pi^2/(6\lambda_3a_n)}}{1-e^{-\pi^2/(2\lambda_3a_n)}},\]

and \(p_{n0*}\mu^0_{a_n}\) is the blocked free photon with the flux
sectors of the torus, which blocking maps to the same sectors. So
stability under iteration, the open clause of Hypothesis P(\(\alpha\))
in the \href{series-parallel-gauge-refinement.md}{series/parallel note},
is reduced for \(U(1)\) in \(D=3\) to exact Gaussian blocking: the
compact part never has to be iterated.

\emph{Proof.} A pushforward does not increase total variation, and
Theorem 2 bounds the distance before the pushforward. The flux of a
coarse 2-torus is the sum of the fluxes of the fine plaquettes it
contains, so the sector label \(m\) is preserved. \(\square\)

The Gaussian comparison that remains, between the blocked free photon
and the free photon at spacing \(a_0\), has total-variation distance of
order one, because the two quadratic forms differ at order one on modes
of the lattice scale, and agree to relative order \(a^2K^2\) on smooth
modes (Proposition 2 and bound (6) of the series/parallel note). It must
be made through smooth observables, where it is explicit because both
sides are Gaussian.

\hypertarget{four-dimensions-the-same-bound-and-why-it-fails}{%
\subsection{3. Four dimensions: the same bound, and why it
fails}\label{four-dimensions-the-same-bound-and-why-it-fails}}

\textbf{Theorem 3.} Let \(D=4\), \(N=L/a\), \(t=g^2\) with \(g^2\le1\).
Then

\[1+12N^4\exp\Bigl\{-\frac{\pi^2}{g^2}-\frac{2\pi^2}{g^2N^4}\Bigr\}\ \le\ R\
\le\ \exp\Bigl\{\frac{8N^4e^{-\pi^2/(8g^2)}}{1-e^{-3\pi^2/(8g^2)}}\Bigr\},\]

\(\|\mu_a-\mu_a^0\|_{\rm TV}\le\log R\), and the flux sectors carry the
weight \(\exp\{-(2\pi^2/g^2)\sum_{\mu<\nu}m_{\mu\nu}^2\}\).

\emph{Proof.} As for Theorem 2, with \(4D=16\), \(|C_3|=4N^4\),
\(|h(m)|^2=|m|^2\), \(6N^4\) plaquettes and two signs for the elementary
loops, loop energy at most \(\frac12\), and harmonic offset \(N^{-4}\).
\(\square\)

\textbf{The recorded failure.} In four dimensions the heat time is the
coupling itself, so at fixed \(g\) the monopole-loop free energy per
lattice cell is bounded above by a fixed number, and the total-variation
bound grows like \((L/a)^4\). The bound is then useless. That the
three-dimensional route itself fails needs an extensive lower bound on
\(\log R\): elementary loops on a sparse sublattice, whose mutual
interaction decays like \(r^{-4}\), give one with a worse constant (a
dilute-gas estimate, expected and not written out here), and an
extensive observable such as \(\sum_p\cos(d\theta)_p\) would then show
that the total-variation distance itself tends to one. The free-energy
density this controls is a vacuum constant. The known continuum
statement is different in kind: the compact theory converges on its
current sector to a renormalized free electromagnetic field
(\href{https://doi.org/10.1007/BF01212424}{Driver 1987}, abstract as
indexed), with the monopole loops renormalizing the charge. A
four-dimensional theorem in the insertion language therefore has to
control observables and a coupling shift per step, as Hypothesis
P(\(\alpha\)) does, rather than the measure. This is the
four-dimensional row of the error budget in the zero-spacing note, now
with an explicit upper bound \(8e^{-\pi^2/(8g^2)}\) on the density per
cell and a single-loop lower bound \(12N^4e^{-\pi^2/g^2}\) on \(R-1\).

\hypertarget{consequence-for-state}{%
\subsection{4. Consequence for STATE}\label{consequence-for-state}}

Item 1 of STATE gains a measure-level theorem for \(U(1)\) in \(D=3\)
with explicit \(a\), \(L\), \(\lambda_3\) (Theorem 2), and a recorded
failure of the same route in \(D=4\) with its reason (Theorem 3). The
non-abelian step that corresponds to Theorem 2 needs a replacement for
the Poisson decomposition: for \(SU(2)\) the exact kernel (11) of the
series/parallel note has image terms of both signs, so Proposition 1 has
no \(SU(2)\) counterpart, and the Laplace-remainder route of its
Proposition 7 remains the path.
\end{document}
